%%%%%%%%%%%%%%%%%%%%%%%%%%%%%%%%%%%%%%%%%%%%%%%%%%%%%%%%%%%%
% 12.8 CONSTRUCTIVE MATHEMATICS
% The Composition of Advanced Mathematics
%%%%%%%%%%%%%%%%%%%%%%%%%%%%%%%%%%%%%%%%%%%%%%%%%%%%%%%%%%%%

\section{Constructive Mathematics}
\label{sec:constructive-mathematics}

\prerequisite{
Formal logic, proof theory, computability theory, recursion theory,
type theory, mathematical induction, real analysis, and familiarity with
classical methods of proof.
}

\learningobjectives{
By the end of this section, the reader should be able to:
\begin{itemize}
    \item explain the central goals of constructive mathematics;
    \item distinguish constructive and classical proof standards;
    \item understand constructive existence;
    \item recognize the computational content of proofs;
    \item understand why witnesses matter;
    \item distinguish direct existence proofs from nonconstructive existence
          proofs;
    \item understand constructive disjunction;
    \item recognize the role of intuitionistic logic;
    \item distinguish classical and intuitionistic reasoning;
    \item understand the law of excluded middle;
    \item understand double-negation elimination;
    \item distinguish \(P\) from \(\neg\neg P\) constructively;
    \item understand proof by contradiction in constructive settings;
    \item recognize proof of negation as a constructive method;
    \item understand constructive implication;
    \item understand constructive universal and existential quantification;
    \item recognize decidable propositions;
    \item distinguish decidable and arbitrary predicates;
    \item understand stability of propositions conceptually;
    \item recognize Markov-style principles conceptually;
    \item understand the constructive meaning of finite search;
    \item recognize countable search limitations;
    \item understand constructive natural numbers;
    \item understand constructive induction and recursion;
    \item recognize algorithm extraction from proofs;
    \item understand constructive real numbers conceptually;
    \item distinguish different constructive representations of real numbers;
    \item understand why equality and order on the reals require care;
    \item recognize locatedness conceptually;
    \item understand apartness relations;
    \item distinguish inequality from apartness;
    \item understand constructive continuity;
    \item recognize constructive intermediate-value issues;
    \item understand constructive compactness conceptually;
    \item recognize total boundedness and completeness;
    \item distinguish constructive compactness principles from classical
          formulations;
    \item understand constructive analysis conceptually;
    \item recognize Bishop-style constructive mathematics;
    \item understand intuitionism conceptually;
    \item recognize finitism and predicativism conceptually;
    \item distinguish constructivism from computability theory;
    \item recognize realizability conceptually;
    \item understand propositions-as-types as a constructive interpretation;
    \item recognize proof interpretations;
    \item understand constructive set theory conceptually;
    \item recognize choice principles in constructive settings;
    \item understand countable and dependent choice conceptually;
    \item recognize the constructive status of the Axiom of Choice;
    \item understand why finite-dimensional mathematics is often highly
          constructive;
    \item recognize constructive algebra;
    \item recognize constructive topology;
    \item understand constructive optimization conceptually;
    \item recognize numerical algorithms as constructive witnesses;
    \item understand error bounds as part of constructive content;
    \item recognize proof-relevant existence in formal verification;
    \item connect constructive mathematics with logic, type theory,
          computability, analysis, numerical mathematics, and foundations.
\end{itemize}
}

%%%%%%%%%%%%%%%%%%%%%%%%%%%%%%%%%%%%%%%%%%%%%%%%%%%%%%%%%%%%
\subsection{What Is Constructive Mathematics?}
\label{subsec:what-is-constructive-mathematics}
%%%%%%%%%%%%%%%%%%%%%%%%%%%%%%%%%%%%%%%%%%%%%%%%%%%%%%%%%%%%

Constructive mathematics is mathematics carried out under proof principles
that emphasize explicit construction and computational meaning.

A classical theorem may establish that an object exists by showing that its
nonexistence would lead to contradiction.

A constructive proof of

\[
\boxed{
\exists x\,P(x)
}
\]

typically requires enough information to actually produce an \(x\) and
verify

\[
\boxed{
P(x).
}
\]

Thus constructive mathematics asks not only whether an object exists, but
also:

\[
\boxed{
\text{How can it be obtained?}
}
\]

%%%%%%%%%%%%%%%%%%%%%%%%%%%%%%%%%%%%%%%%%%%%%%%%%%%%%%%%%%%%
\subsection{Constructive Existence}
\label{subsec:constructive-existence}
%%%%%%%%%%%%%%%%%%%%%%%%%%%%%%%%%%%%%%%%%%%%%%%%%%%%%%%%%%%%

Under a constructive interpretation, a proof of

\[
\exists x\in A\,P(x)
\]

contains two pieces of information:

\[
\boxed{
a\in A
}
\]

and

\[
\boxed{
\text{a proof of }P(a).
}
\]

Thus existential proof behaves like a pair

\[
\boxed{
(a,p).
}
\]

This matches the dependent-pair interpretation developed in
Section~\ref{sec:type-theory}.

%%%%%%%%%%%%%%%%%%%%%%%%%%%%%%%%%%%%%%%%%%%%%%%%%%%%%%%%%%%%
\subsection{Worked Example: Constructive Existence}
\label{subsec:worked-constructive-existence}
%%%%%%%%%%%%%%%%%%%%%%%%%%%%%%%%%%%%%%%%%%%%%%%%%%%%%%%%%%%%

\begin{workedexamplebox}{An Explicit Even Prime}

Consider the statement

\[
\exists p\in\mathbb{N}
\,
(
p\text{ is prime and even}
).
\]

Choose

\[
p=2.
\]

The integer \(2\) is divisible only by \(1\) and itself, and

\[
2=2\cdot1.
\]

Hence \(2\) is both prime and even.

\begin{answerbox}
\[
\boxed{
p=2
}
\]

is an explicit witness.
\end{answerbox}

The proof is constructive because it exhibits the required object.

\end{workedexamplebox}

%%%%%%%%%%%%%%%%%%%%%%%%%%%%%%%%%%%%%%%%%%%%%%%%%%%%%%%%%%%%
\subsection{Nonconstructive Existence}
\label{subsec:nonconstructive-existence}
%%%%%%%%%%%%%%%%%%%%%%%%%%%%%%%%%%%%%%%%%%%%%%%%%%%%%%%%%%%%

A proof is nonconstructive when it establishes existence without giving a
method to construct a witness.

A common classical pattern is:

\[
\boxed{
\neg
\left(
\forall x\,\neg P(x)
\right)
}
\]

and then conclude

\[
\boxed{
\exists x\,P(x).
}
\]

Classically this is valid.

Constructively, the step from

\[
\neg\forall x\,\neg P(x)
\]

to

\[
\exists x\,P(x)
\]

is not generally available.

%%%%%%%%%%%%%%%%%%%%%%%%%%%%%%%%%%%%%%%%%%%%%%%%%%%%%%%%%%%%
\subsection{Classical Versus Constructive Logic}
\label{subsec:classical-versus-constructive-logic}
%%%%%%%%%%%%%%%%%%%%%%%%%%%%%%%%%%%%%%%%%%%%%%%%%%%%%%%%%%%%

Constructive mathematics is often formalized using intuitionistic logic.

Classical logic accepts principles such as

\[
\boxed{
P\lor\neg P
}
\]

for every proposition \(P\).

This is the law of excluded middle.

Constructive logic does not accept this as an unrestricted theorem.

%%%%%%%%%%%%%%%%%%%%%%%%%%%%%%%%%%%%%%%%%%%%%%%%%%%%%%%%%%%%
\subsection{Law of Excluded Middle}
\label{subsec:constructive-excluded-middle}
%%%%%%%%%%%%%%%%%%%%%%%%%%%%%%%%%%%%%%%%%%%%%%%%%%%%%%%%%%%%

The law of excluded middle is

\[
\boxed{
P\lor\neg P.
}
\]

Constructively, proving a disjunction means providing evidence for one side.

Therefore a proof of

\[
P\lor\neg P
\]

would require a method to determine whether \(P\) holds.

For arbitrary propositions, such a method may not exist.

%%%%%%%%%%%%%%%%%%%%%%%%%%%%%%%%%%%%%%%%%%%%%%%%%%%%%%%%%%%%
\subsection{Decidable Propositions}
\label{subsec:decidable-propositions}
%%%%%%%%%%%%%%%%%%%%%%%%%%%%%%%%%%%%%%%%%%%%%%%%%%%%%%%%%%%%

A proposition \(P\) is decidable if one can constructively prove

\[
\boxed{
P\lor\neg P.
}
\]

Thus excluded middle remains valid for propositions for which a decision
procedure is available.

For example, for concrete integers \(m,n\),

\[
\boxed{
m=n
\lor
m\neq n
}
\]

is decidable.

%%%%%%%%%%%%%%%%%%%%%%%%%%%%%%%%%%%%%%%%%%%%%%%%%%%%%%%%%%%%
\subsection{Finite Search Is Constructive}
\label{subsec:finite-search-constructive}
%%%%%%%%%%%%%%%%%%%%%%%%%%%%%%%%%%%%%%%%%%%%%%%%%%%%%%%%%%%%

Suppose

\[
A=\{a_1,\ldots,a_n\}
\]

is finite and property \(P(a_i)\) is decidable for each \(a_i\).

Then one may examine each element.

Hence

\[
\boxed{
\exists a\in A\,P(a)
}
\]

is decidable by finite search.

Similarly,

\[
\boxed{
\forall a\in A\,P(a)
}
\]

can be checked constructively.

%%%%%%%%%%%%%%%%%%%%%%%%%%%%%%%%%%%%%%%%%%%%%%%%%%%%%%%%%%%%
\subsection{Worked Example: Finite Constructive Search}
\label{subsec:worked-finite-constructive-search}
%%%%%%%%%%%%%%%%%%%%%%%%%%%%%%%%%%%%%%%%%%%%%%%%%%%%%%%%%%%%

\begin{workedexamplebox}{Finding a Multiple of \(3\)}

Let

\[
A=\{4,7,9,10\}.
\]

Check divisibility by \(3\):

\[
4\not\equiv0\pmod3,
\]

\[
7\not\equiv0\pmod3,
\]

\[
9\equiv0\pmod3.
\]

Therefore,

\begin{answerbox}
\[
\boxed{
9
}
\]

is a constructive witness to

\[
\exists a\in A
\,
(3\mid a).
\]
\end{answerbox}

\end{workedexamplebox}

%%%%%%%%%%%%%%%%%%%%%%%%%%%%%%%%%%%%%%%%%%%%%%%%%%%%%%%%%%%%
\subsection{Double Negation}
\label{subsec:constructive-double-negation}
%%%%%%%%%%%%%%%%%%%%%%%%%%%%%%%%%%%%%%%%%%%%%%%%%%%%%%%%%%%%

Constructively, one can prove

\[
\boxed{
P
\rightarrow
\neg\neg P.
}
\]

If \(P\) holds, then assuming \(\neg P\) gives contradiction.

However, the converse

\[
\boxed{
\neg\neg P
\rightarrow
P
}
\]

is not generally valid intuitionistically.

Thus \(P\) and \(\neg\neg P\) must be distinguished.

%%%%%%%%%%%%%%%%%%%%%%%%%%%%%%%%%%%%%%%%%%%%%%%%%%%%%%%%%%%%
\subsection{Proof by Contradiction}
\label{subsec:constructive-proof-by-contradiction}
%%%%%%%%%%%%%%%%%%%%%%%%%%%%%%%%%%%%%%%%%%%%%%%%%%%%%%%%%%%%

Classically, to prove \(P\), one may assume

\[
\neg P
\]

and derive contradiction.

This proves

\[
\neg\neg P.
\]

Classical double-negation elimination then gives \(P\).

Constructively, the final step is unavailable in general.

Therefore proof by contradiction is not universally accepted as a method
for proving arbitrary positive statements.

%%%%%%%%%%%%%%%%%%%%%%%%%%%%%%%%%%%%%%%%%%%%%%%%%%%%%%%%%%%%
\subsection{Constructive Proof of Negation}
\label{subsec:constructive-proof-negation}
%%%%%%%%%%%%%%%%%%%%%%%%%%%%%%%%%%%%%%%%%%%%%%%%%%%%%%%%%%%%

Proof by contradiction remains fully constructive when the goal itself is a
negation.

To prove

\[
\boxed{
\neg P,
}
\]

assume

\[
P
\]

and derive

\[
\bot.
\]

Since

\[
\boxed{
\neg P
=
P\rightarrow\bot,
}
\]

this directly constructs the required proof.

%%%%%%%%%%%%%%%%%%%%%%%%%%%%%%%%%%%%%%%%%%%%%%%%%%%%%%%%%%%%
\subsection{Worked Example: Constructive Negation}
\label{subsec:worked-constructive-negation}
%%%%%%%%%%%%%%%%%%%%%%%%%%%%%%%%%%%%%%%%%%%%%%%%%%%%%%%%%%%%

\begin{workedexamplebox}{No Integer Lies Strictly Between \(2\) and \(3\)}

Suppose

\[
n\in\mathbb{Z}
\]

satisfies

\[
2<n<3.
\]

Since \(n\) is an integer and \(n>2\),

\[
n\geq3.
\]

This contradicts

\[
n<3.
\]

Therefore,

\begin{answerbox}
\[
\boxed{
\neg
\exists n\in\mathbb{Z}
\,
(2<n<3).
}
\]
\end{answerbox}

This contradiction proof is constructive because the target is a negation.

\end{workedexamplebox}

%%%%%%%%%%%%%%%%%%%%%%%%%%%%%%%%%%%%%%%%%%%%%%%%%%%%%%%%%%%%
\subsection{Constructive Disjunction}
\label{subsec:constructive-disjunction}
%%%%%%%%%%%%%%%%%%%%%%%%%%%%%%%%%%%%%%%%%%%%%%%%%%%%%%%%%%%%

To prove

\[
P\lor Q
\]

constructively, one must establish either:

\[
\boxed{
P
}
\]

or

\[
\boxed{
Q.
}
\]

A proof merely showing that

\[
\neg(\neg P\land\neg Q)
\]

does not necessarily determine which disjunct holds.

Thus constructive disjunction contains information.

%%%%%%%%%%%%%%%%%%%%%%%%%%%%%%%%%%%%%%%%%%%%%%%%%%%%%%%%%%%%
\subsection{Constructive Implication}
\label{subsec:constructive-implication}
%%%%%%%%%%%%%%%%%%%%%%%%%%%%%%%%%%%%%%%%%%%%%%%%%%%%%%%%%%%%

A constructive proof of

\[
P\rightarrow Q
\]

is a method transforming any proof of \(P\) into a proof of \(Q\).

Under Curry--Howard, this is represented by a function

\[
\boxed{
P\to Q.
}
\]

Implication therefore has explicit operational meaning.

%%%%%%%%%%%%%%%%%%%%%%%%%%%%%%%%%%%%%%%%%%%%%%%%%%%%%%%%%%%%
\subsection{Constructive Universal Quantification}
\label{subsec:constructive-universal}
%%%%%%%%%%%%%%%%%%%%%%%%%%%%%%%%%%%%%%%%%%%%%%%%%%%%%%%%%%%%

A proof of

\[
\forall x\in A\,P(x)
\]

constructively gives a uniform method that takes arbitrary

\[
a\in A
\]

and produces a proof of

\[
\boxed{
P(a).
}
\]

Thus universal quantification corresponds to a function assigning proofs
uniformly over the domain.

%%%%%%%%%%%%%%%%%%%%%%%%%%%%%%%%%%%%%%%%%%%%%%%%%%%%%%%%%%%%
\subsection{Constructive Existential Quantification}
\label{subsec:constructive-existential}
%%%%%%%%%%%%%%%%%%%%%%%%%%%%%%%%%%%%%%%%%%%%%%%%%%%%%%%%%%%%

A constructive proof of

\[
\exists x\in A\,P(x)
\]

must provide:

\[
\boxed{
\text{a witness }a\in A
}
\]

and

\[
\boxed{
\text{verification of }P(a).
}
\]

This is stronger than merely proving that assuming no witness exists leads
to contradiction.

%%%%%%%%%%%%%%%%%%%%%%%%%%%%%%%%%%%%%%%%%%%%%%%%%%%%%%%%%%%%
\subsection{Brouwer--Heyting--Kolmogorov Interpretation}
\label{subsec:bhk-interpretation}
%%%%%%%%%%%%%%%%%%%%%%%%%%%%%%%%%%%%%%%%%%%%%%%%%%%%%%%%%%%%

The constructive meaning of logical connectives is often summarized by the
Brouwer--Heyting--Kolmogorov interpretation.

Informally:

\[
\boxed{
P\land Q
}
\]

means a construction of both \(P\) and \(Q\);

\[
\boxed{
P\lor Q
}
\]

means a construction of one side together with information identifying it;

\[
\boxed{
P\rightarrow Q
}
\]

means a construction transforming proofs of \(P\) into proofs of \(Q\);

\[
\boxed{
\exists x\,P(x)
}
\]

means a witness and a proof;

\[
\boxed{
\forall x\,P(x)
}
\]

means a uniform construction producing a proof for each \(x\).

%%%%%%%%%%%%%%%%%%%%%%%%%%%%%%%%%%%%%%%%%%%%%%%%%%%%%%%%%%%%
\subsection{Constructive Truth}
\label{subsec:constructive-truth}
%%%%%%%%%%%%%%%%%%%%%%%%%%%%%%%%%%%%%%%%%%%%%%%%%%%%%%%%%%%%

In a constructive setting, truth is often associated with possession of
evidence or a construction.

This does not mean mathematical truth is reduced to physical calculation.

Rather, the formal meaning of a proposition is tied to what would count as
a proof of it.

%%%%%%%%%%%%%%%%%%%%%%%%%%%%%%%%%%%%%%%%%%%%%%%%%%%%%%%%%%%%
\subsection{Stability}
\label{subsec:constructive-stability}
%%%%%%%%%%%%%%%%%%%%%%%%%%%%%%%%%%%%%%%%%%%%%%%%%%%%%%%%%%%%

A proposition \(P\) is called stable when

\[
\boxed{
\neg\neg P
\rightarrow
P
}
\]

is constructively provable.

Some propositions are stable even though unrestricted double-negation
elimination is not accepted.

For example, many negative propositions have this property.

%%%%%%%%%%%%%%%%%%%%%%%%%%%%%%%%%%%%%%%%%%%%%%%%%%%%%%%%%%%%
\subsection{Negative Formulas}
\label{subsec:negative-formulas-constructive}
%%%%%%%%%%%%%%%%%%%%%%%%%%%%%%%%%%%%%%%%%%%%%%%%%%%%%%%%%%%%

Statements already expressed primarily through negation can behave more
classically than arbitrary positive statements.

For instance, a statement of the form

\[
\boxed{
\neg Q
}
\]

is stable because

\[
\neg\neg\neg Q
\rightarrow
\neg Q
\]

is intuitionistically valid.

Thus constructive reasoning distinguishes logical form carefully.

%%%%%%%%%%%%%%%%%%%%%%%%%%%%%%%%%%%%%%%%%%%%%%%%%%%%%%%%%%%%
\subsection{Markov-Type Principles}
\label{subsec:markov-principle-preview}
%%%%%%%%%%%%%%%%%%%%%%%%%%%%%%%%%%%%%%%%%%%%%%%%%%%%%%%%%%%%

For decidable predicates on natural numbers, one may consider principles of
the form

\[
\boxed{
\neg\neg
\exists n\,P(n)
\rightarrow
\exists n\,P(n).
}
\]

Such principles are not generally theorems of intuitionistic logic alone.

Their acceptance depends on the particular constructive foundation being
used.

This illustrates that ``constructive mathematics'' is not one unique formal
system.

%%%%%%%%%%%%%%%%%%%%%%%%%%%%%%%%%%%%%%%%%%%%%%%%%%%%%%%%%%%%
\subsection{Varieties of Constructive Mathematics}
\label{subsec:varieties-constructive-mathematics}
%%%%%%%%%%%%%%%%%%%%%%%%%%%%%%%%%%%%%%%%%%%%%%%%%%%%%%%%%%%%

Constructive mathematics includes several related approaches.

Examples include:

\[
\boxed{
\text{intuitionistic mathematics},
}
\]

\[
\boxed{
\text{Bishop-style constructive mathematics},
}
\]

\[
\boxed{
\text{recursive or computable mathematics},
}
\]

and

\[
\boxed{
\text{type-theoretic constructive mathematics}.
}
\]

They share many themes but do not accept exactly the same principles.

%%%%%%%%%%%%%%%%%%%%%%%%%%%%%%%%%%%%%%%%%%%%%%%%%%%%%%%%%%%%
\subsection{Bishop-Style Constructive Mathematics}
\label{subsec:bishop-constructive-mathematics}
%%%%%%%%%%%%%%%%%%%%%%%%%%%%%%%%%%%%%%%%%%%%%%%%%%%%%%%%%%%%

Bishop-style constructive mathematics aims to develop substantial ordinary
mathematics using methods with explicit computational meaning.

A major goal is to prove theorems that remain valid under several different
constructive interpretations.

This style often emphasizes:

\[
\boxed{
\text{explicit estimates},
}
\]

\[
\boxed{
\text{algorithms},
}
\]

and

\[
\boxed{
\text{quantitative information}.
}
\]

%%%%%%%%%%%%%%%%%%%%%%%%%%%%%%%%%%%%%%%%%%%%%%%%%%%%%%%%%%%%
\subsection{Constructive Mathematics Is Not Merely Computable Mathematics}
\label{subsec:constructive-vs-computable}
%%%%%%%%%%%%%%%%%%%%%%%%%%%%%%%%%%%%%%%%%%%%%%%%%%%%%%%%%%%%

Constructive mathematics and computability theory are strongly connected
but distinct.

Computability theory studies what algorithms can compute.

Constructive mathematics studies what forms of proof are mathematically
acceptable under constructive interpretations.

A constructive proof often yields computational content, but the precise
relationship depends on the formal system.

%%%%%%%%%%%%%%%%%%%%%%%%%%%%%%%%%%%%%%%%%%%%%%%%%%%%%%%%%%%%
\subsection{Constructive Natural Numbers}
\label{subsec:constructive-natural-numbers}
%%%%%%%%%%%%%%%%%%%%%%%%%%%%%%%%%%%%%%%%%%%%%%%%%%%%%%%%%%%%

Natural numbers are especially well suited to constructive mathematics.

They are generated inductively by

\[
\boxed{
0
}
\]

and

\[
\boxed{
S(n)=n+1.
}
\]

Equality of natural numbers is decidable.

Finite calculations and induction can therefore be carried out
constructively in a direct way.

%%%%%%%%%%%%%%%%%%%%%%%%%%%%%%%%%%%%%%%%%%%%%%%%%%%%%%%%%%%%
\subsection{Constructive Induction}
\label{subsec:constructive-induction}
%%%%%%%%%%%%%%%%%%%%%%%%%%%%%%%%%%%%%%%%%%%%%%%%%%%%%%%%%%%%

Mathematical induction is constructive.

To prove

\[
\forall n\in\mathbb{N}\,P(n),
\]

one constructs:

\[
\boxed{
P(0)
}
\]

and a method

\[
\boxed{
P(n)
\rightarrow
P(n+1).
}
\]

The induction principle then produces evidence for every natural number.

%%%%%%%%%%%%%%%%%%%%%%%%%%%%%%%%%%%%%%%%%%%%%%%%%%%%%%%%%%%%
\subsection{Constructive Recursion}
\label{subsec:constructive-recursion}
%%%%%%%%%%%%%%%%%%%%%%%%%%%%%%%%%%%%%%%%%%%%%%%%%%%%%%%%%%%%

Recursive definitions on the natural numbers provide explicit algorithms.

For example,

\[
f(0)=a,
\]

\[
f(n+1)=G(n,f(n))
\]

gives a direct procedure for calculating \(f(n)\) from \(n\).

Thus recursion is naturally constructive.

%%%%%%%%%%%%%%%%%%%%%%%%%%%%%%%%%%%%%%%%%%%%%%%%%%%%%%%%%%%%
\subsection{Algorithm Extraction}
\label{subsec:constructive-algorithm-extraction}
%%%%%%%%%%%%%%%%%%%%%%%%%%%%%%%%%%%%%%%%%%%%%%%%%%%%%%%%%%%%

Suppose one proves constructively

\[
\boxed{
\forall x\in A
\exists y\in B
\,
R(x,y).
}
\]

The proof may contain enough information to produce a function

\[
\boxed{
f:A\to B
}
\]

such that

\[
\boxed{
R(x,f(x))
}
\]

for every \(x\).

This is algorithm extraction from proof.

%%%%%%%%%%%%%%%%%%%%%%%%%%%%%%%%%%%%%%%%%%%%%%%%%%%%%%%%%%%%
\subsection{Worked Example: Extracting a Witness Function}
\label{subsec:worked-witness-function}
%%%%%%%%%%%%%%%%%%%%%%%%%%%%%%%%%%%%%%%%%%%%%%%%%%%%%%%%%%%%

\begin{workedexamplebox}{A Larger Natural Number}

Consider

\[
\forall n\in\mathbb{N}
\exists m\in\mathbb{N}
\,
(m>n).
\]

Given \(n\), choose

\[
m=n+1.
\]

Since

\[
n+1>n,
\]

the witness is valid.

The proof directly yields the function

\[
\begin{answerbox}
\[
\boxed{
f(n)=n+1.
}
\]
\end{answerbox}

\end{workedexamplebox}

%%%%%%%%%%%%%%%%%%%%%%%%%%%%%%%%%%%%%%%%%%%%%%%%%%%%%%%%%%%%
\subsection{Constructive Real Numbers}
\label{subsec:constructive-real-numbers}
%%%%%%%%%%%%%%%%%%%%%%%%%%%%%%%%%%%%%%%%%%%%%%%%%%%%%%%%%%%%

Real numbers require more care constructively than natural numbers.

A real number may be represented by effective approximation data, such as a
Cauchy sequence of rational numbers together with a modulus controlling its
convergence.

The representation must contain enough information to approximate the real
to any desired precision.

%%%%%%%%%%%%%%%%%%%%%%%%%%%%%%%%%%%%%%%%%%%%%%%%%%%%%%%%%%%%
\subsection{Cauchy Representation}
\label{subsec:constructive-cauchy-representation}
%%%%%%%%%%%%%%%%%%%%%%%%%%%%%%%%%%%%%%%%%%%%%%%%%%%%%%%%%%%%

A constructive real number may be represented by rational approximations

\[
q_0,q_1,q_2,\ldots
\]

with explicit control such as

\[
\boxed{
|q_m-q_n|
<
2^{-k}
}
\]

whenever

\[
m,n
\]

are beyond a computably specified stage depending on \(k\).

The convergence information is part of the data.

%%%%%%%%%%%%%%%%%%%%%%%%%%%%%%%%%%%%%%%%%%%%%%%%%%%%%%%%%%%%
\subsection{Modulus of Convergence}
\label{subsec:modulus-of-convergence}
%%%%%%%%%%%%%%%%%%%%%%%%%%%%%%%%%%%%%%%%%%%%%%%%%%%%%%%%%%%%

A modulus of convergence is a function

\[
N:\mathbb{N}\to\mathbb{N}
\]

such that for sufficiently large indices,

\[
m,n\geq N(k)
\]

implies

\[
\boxed{
|q_m-q_n|
<
2^{-k}.
}
\]

Thus the proof of Cauchy convergence contains a quantitative rate.

%%%%%%%%%%%%%%%%%%%%%%%%%%%%%%%%%%%%%%%%%%%%%%%%%%%%%%%%%%%%
\subsection{Why Rates Matter}
\label{subsec:why-rates-matter-constructive}
%%%%%%%%%%%%%%%%%%%%%%%%%%%%%%%%%%%%%%%%%%%%%%%%%%%%%%%%%%%%

Classically, knowing

\[
q_n\to x
\]

may be enough.

Constructively, one often needs to know how far to compute to guarantee an
error less than a specified tolerance.

Thus the statement

\[
\boxed{
|q_n-x|<\varepsilon
}
\]

should come with a method for determining a sufficient \(n\).

%%%%%%%%%%%%%%%%%%%%%%%%%%%%%%%%%%%%%%%%%%%%%%%%%%%%%%%%%%%%
\subsection{Worked Example: Explicit Approximation}
\label{subsec:worked-explicit-approximation}
%%%%%%%%%%%%%%%%%%%%%%%%%%%%%%%%%%%%%%%%%%%%%%%%%%%%%%%%%%%%

\begin{workedexamplebox}{Approximating a Limit}

Suppose

\[
q_n
=
1+\frac{1}{n+1}.
\]

Then

\[
q_n\to1.
\]

To ensure

\[
|q_n-1|
<
\varepsilon,
\]

we need

\[
\frac{1}{n+1}<\varepsilon.
\]

It is sufficient to choose

\[
n+1>\frac1\varepsilon.
\]

Thus one explicit choice is

\[
\begin{answerbox}
\[
\boxed{
n
\geq
\left\lceil\frac1\varepsilon\right\rceil.
}
\]
\end{answerbox}

The proof supplies a quantitative approximation procedure.

\end{workedexamplebox}

%%%%%%%%%%%%%%%%%%%%%%%%%%%%%%%%%%%%%%%%%%%%%%%%%%%%%%%%%%%%
\subsection{Equality of Constructive Reals}
\label{subsec:constructive-real-equality}
%%%%%%%%%%%%%%%%%%%%%%%%%%%%%%%%%%%%%%%%%%%%%%%%%%%%%%%%%%%%

Equality of arbitrary real numbers is not generally decidable
constructively.

Given two real numbers \(x\) and \(y\), finite approximations may not reveal
whether

\[
\boxed{
x=y
}
\]

or merely

\[
\boxed{
x
\text{ and }
y
\text{ are extremely close}.
}
\]

Therefore constructive real analysis often uses stronger positive notions
such as apartness.

%%%%%%%%%%%%%%%%%%%%%%%%%%%%%%%%%%%%%%%%%%%%%%%%%%%%%%%%%%%%
\subsection{Apartness}
\label{subsec:apartness}
%%%%%%%%%%%%%%%%%%%%%%%%%%%%%%%%%%%%%%%%%%%%%%%%%%%%%%%%%%%%

An apartness relation

\[
\boxed{
x\# y
}
\]

expresses positive evidence that \(x\) and \(y\) are distinct.

For real numbers, one may interpret

\[
x\# y
\]

as having explicit evidence that

\[
\boxed{
|x-y|>\varepsilon
}
\]

for some positive rational \(\varepsilon\).

This is stronger than merely

\[
\neg(x=y).
\]

%%%%%%%%%%%%%%%%%%%%%%%%%%%%%%%%%%%%%%%%%%%%%%%%%%%%%%%%%%%%
\subsection{Inequality Versus Apartness}
\label{subsec:inequality-versus-apartness}
%%%%%%%%%%%%%%%%%%%%%%%%%%%%%%%%%%%%%%%%%%%%%%%%%%%%%%%%%%%%

Constructively,

\[
\boxed{
x\# y
\Longrightarrow
x\neq y.
}
\]

The converse need not be derivable.

The statement

\[
x\neq y
\]

only means equality leads to contradiction.

Apartness provides positive quantitative separation.

%%%%%%%%%%%%%%%%%%%%%%%%%%%%%%%%%%%%%%%%%%%%%%%%%%%%%%%%%%%%
\subsection{Positive Order}
\label{subsec:constructive-positive-order}
%%%%%%%%%%%%%%%%%%%%%%%%%%%%%%%%%%%%%%%%%%%%%%%%%%%%%%%%%%%%

For constructive real numbers, a relation such as

\[
x<y
\]

is often interpreted positively.

One seeks explicit rational separation showing that \(y-x\) is bounded
away from zero.

This makes strict order computationally meaningful.

%%%%%%%%%%%%%%%%%%%%%%%%%%%%%%%%%%%%%%%%%%%%%%%%%%%%%%%%%%%%
\subsection{Locatedness}
\label{subsec:locatedness}
%%%%%%%%%%%%%%%%%%%%%%%%%%%%%%%%%%%%%%%%%%%%%%%%%%%%%%%%%%%%

Locatedness expresses the ability to approximate the position of a point or
set relative to rational bounds.

For a subset \(A\) of a metric space, one may seek constructive information
about

\[
\boxed{
d(x,A)
=
\inf_{a\in A}d(x,a).
}
\]

A located set is one for which this distance can be meaningfully
approximated within the constructive framework.

%%%%%%%%%%%%%%%%%%%%%%%%%%%%%%%%%%%%%%%%%%%%%%%%%%%%%%%%%%%%
\subsection{Constructive Suprema}
\label{subsec:constructive-suprema}
%%%%%%%%%%%%%%%%%%%%%%%%%%%%%%%%%%%%%%%%%%%%%%%%%%%%%%%%%%%%

Classically, every nonempty bounded-above subset of \(\mathbb{R}\) has a
supremum.

Constructively, arbitrary subsets may not provide enough information to
compute or locate such a supremum.

Additional hypotheses are often needed.

For explicitly given sequences or located sets, constructive supremum
results can frequently be obtained.

%%%%%%%%%%%%%%%%%%%%%%%%%%%%%%%%%%%%%%%%%%%%%%%%%%%%%%%%%%%%
\subsection{Constructive Continuity}
\label{subsec:constructive-continuity}
%%%%%%%%%%%%%%%%%%%%%%%%%%%%%%%%%%%%%%%%%%%%%%%%%%%%%%%%%%%%

Continuity fits naturally with constructive approximation.

For

\[
f:X\to Y,
\]

one seeks a method turning desired output precision into sufficient input
precision.

For metric spaces, this resembles the usual \(\varepsilon\)-\(\delta\)
definition, but the dependence should be computationally meaningful.

%%%%%%%%%%%%%%%%%%%%%%%%%%%%%%%%%%%%%%%%%%%%%%%%%%%%%%%%%%%%
\subsection{Uniform Continuity}
\label{subsec:constructive-uniform-continuity}
%%%%%%%%%%%%%%%%%%%%%%%%%%%%%%%%%%%%%%%%%%%%%%%%%%%%%%%%%%%%

Uniform continuity provides particularly strong constructive information.

For each

\[
\varepsilon>0,
\]

one finds a single

\[
\delta>0
\]

such that

\[
d_X(x,y)<\delta
\]

implies

\[
\boxed{
d_Y(f(x),f(y))<\varepsilon
}
\]

for all \(x,y\) in the domain.

A modulus of uniform continuity can encode this dependence explicitly.

%%%%%%%%%%%%%%%%%%%%%%%%%%%%%%%%%%%%%%%%%%%%%%%%%%%%%%%%%%%%
\subsection{Modulus of Continuity}
\label{subsec:modulus-of-continuity}
%%%%%%%%%%%%%%%%%%%%%%%%%%%%%%%%%%%%%%%%%%%%%%%%%%%%%%%%%%%%

A modulus of continuity gives a quantitative rule for choosing input
precision.

For example, one may have a function

\[
\omega:\mathbb{N}\to\mathbb{N}
\]

such that

\[
|x-y|<2^{-\omega(k)}
\]

implies

\[
\boxed{
|f(x)-f(y)|<2^{-k}.
}
\]

This makes continuity algorithmically usable.

%%%%%%%%%%%%%%%%%%%%%%%%%%%%%%%%%%%%%%%%%%%%%%%%%%%%%%%%%%%%
\subsection{Constructive Intermediate Value Issues}
\label{subsec:constructive-intermediate-value}
%%%%%%%%%%%%%%%%%%%%%%%%%%%%%%%%%%%%%%%%%%%%%%%%%%%%%%%%%%%%

Classically, if continuous \(f\) satisfies

\[
f(a)<0<f(b),
\]

then there exists

\[
c\in(a,b)
\]

with

\[
f(c)=0.
\]

Constructively, exact root existence can require additional care because
equality to zero may not be decidable.

Approximate root statements are often more directly constructive.

%%%%%%%%%%%%%%%%%%%%%%%%%%%%%%%%%%%%%%%%%%%%%%%%%%%%%%%%%%%%
\subsection{Approximate Intermediate Value Principle}
\label{subsec:approximate-intermediate-value}
%%%%%%%%%%%%%%%%%%%%%%%%%%%%%%%%%%%%%%%%%%%%%%%%%%%%%%%%%%%%

A constructive version may seek, for every

\[
\varepsilon>0,
\]

a point \(c\) satisfying

\[
\boxed{
|f(c)|<\varepsilon.
}
\]

Such a point can often be produced by a bisection-style algorithm when
sufficient sign information is available.

%%%%%%%%%%%%%%%%%%%%%%%%%%%%%%%%%%%%%%%%%%%%%%%%%%%%%%%%%%%%
\subsection{Worked Example: Constructive Bisection}
\label{subsec:worked-constructive-bisection}
%%%%%%%%%%%%%%%%%%%%%%%%%%%%%%%%%%%%%%%%%%%%%%%%%%%%%%%%%%%%

\begin{workedexamplebox}{Approximating a Root}

Suppose \(f\) is continuous on \([a,b]\) and one has verified

\[
f(a)<0
\]

and

\[
f(b)>0.
\]

Let

\[
m=\frac{a+b}{2}.
\]

Evaluate enough information about \(f(m)\) to choose a smaller interval in
which the sign change is retained.

After \(n\) bisections, the interval width is

\[
\frac{b-a}{2^n}.
\]

To obtain width below \(\varepsilon\), choose

\[
n
>
\log_2
\left(
\frac{b-a}{\varepsilon}
\right).
\]

\begin{answerbox}
\[
\boxed{
\text{Bisection produces increasingly accurate root-location data.}
}
\]
\end{answerbox}

\end{workedexamplebox}

%%%%%%%%%%%%%%%%%%%%%%%%%%%%%%%%%%%%%%%%%%%%%%%%%%%%%%%%%%%%
\subsection{Constructive Compactness}
\label{subsec:constructive-compactness}
%%%%%%%%%%%%%%%%%%%%%%%%%%%%%%%%%%%%%%%%%%%%%%%%%%%%%%%%%%%%

Classical topology contains many equivalent definitions of compactness.

Constructively, these equivalences can separate.

For metric spaces, a useful constructive combination is often:

\[
\boxed{
\text{complete}
+
\text{totally bounded}.
}
\]

Total boundedness provides finite approximations at every scale.

%%%%%%%%%%%%%%%%%%%%%%%%%%%%%%%%%%%%%%%%%%%%%%%%%%%%%%%%%%%%
\subsection{Total Boundedness}
\label{subsec:constructive-total-boundedness}
%%%%%%%%%%%%%%%%%%%%%%%%%%%%%%%%%%%%%%%%%%%%%%%%%%%%%%%%%%%%

A metric space \(X\) is totally bounded if for every

\[
\varepsilon>0,
\]

one can cover \(X\) by finitely many balls of radius \(\varepsilon\).

Constructively, this means being able to produce such finite
\(\varepsilon\)-nets.

Thus total boundedness contains explicit approximation information.

%%%%%%%%%%%%%%%%%%%%%%%%%%%%%%%%%%%%%%%%%%%%%%%%%%%%%%%%%%%%
\subsection{Constructive Completeness}
\label{subsec:constructive-completeness}
%%%%%%%%%%%%%%%%%%%%%%%%%%%%%%%%%%%%%%%%%%%%%%%%%%%%%%%%%%%%

A metric space is complete when Cauchy sequences converge.

Constructively, the Cauchy sequence should usually carry enough convergence
information to determine approximations effectively.

Thus completeness is tied to explicit Cauchy data rather than merely an
existential limit statement.

%%%%%%%%%%%%%%%%%%%%%%%%%%%%%%%%%%%%%%%%%%%%%%%%%%%%%%%%%%%%
\subsection{Finite Nets}
\label{subsec:finite-nets}
%%%%%%%%%%%%%%%%%%%%%%%%%%%%%%%%%%%%%%%%%%%%%%%%%%%%%%%%%%%%

An \(\varepsilon\)-net is a finite set

\[
\boxed{
\{x_1,\ldots,x_n\}
}
\]

such that every point \(x\in X\) satisfies

\[
\boxed{
d(x,x_i)<\varepsilon
}
\]

for some \(i\).

Producing such nets at arbitrary precision makes compact spaces
computationally approximable.

%%%%%%%%%%%%%%%%%%%%%%%%%%%%%%%%%%%%%%%%%%%%%%%%%%%%%%%%%%%%
\subsection{Constructive Extreme Value Principles}
\label{subsec:constructive-extreme-value}
%%%%%%%%%%%%%%%%%%%%%%%%%%%%%%%%%%%%%%%%%%%%%%%%%%%%%%%%%%%%

For a uniformly continuous function on a totally bounded space, one can
often approximate the supremum and infimum constructively.

Exact attainment at a point may require additional hypotheses or a
carefully formulated theorem.

This distinction separates:

\[
\boxed{
\text{approximating an optimal value}
}
\]

from

\[
\boxed{
\text{constructing an exact optimizer}.
}
\]

%%%%%%%%%%%%%%%%%%%%%%%%%%%%%%%%%%%%%%%%%%%%%%%%%%%%%%%%%%%%
\subsection{Approximate Optimization}
\label{subsec:constructive-approximate-optimization}
%%%%%%%%%%%%%%%%%%%%%%%%%%%%%%%%%%%%%%%%%%%%%%%%%%%%%%%%%%%%

Suppose

\[
f:X\to\mathbb{R}
\]

is uniformly continuous and \(X\) can be approximated by finite nets.

For each tolerance

\[
\varepsilon>0,
\]

one may search a sufficiently fine finite net and construct \(x_\varepsilon\)
such that

\[
\boxed{
f(x_\varepsilon)
\leq
\inf_{x\in X}f(x)
+
\varepsilon.
}
\]

This is a naturally constructive optimization statement.

%%%%%%%%%%%%%%%%%%%%%%%%%%%%%%%%%%%%%%%%%%%%%%%%%%%%%%%%%%%%
\subsection{Exact Versus Approximate Solutions}
\label{subsec:exact-versus-approximate-constructive}
%%%%%%%%%%%%%%%%%%%%%%%%%%%%%%%%%%%%%%%%%%%%%%%%%%%%%%%%%%%%

Constructive mathematics frequently replaces an exact existential statement
with controlled approximation.

Instead of merely claiming

\[
\boxed{
\exists x^*
\,
f(x^*)=m,
}
\]

one may produce, for every \(\varepsilon>0\),

\[
\boxed{
x_\varepsilon
}
\]

satisfying

\[
\boxed{
f(x_\varepsilon)
<
m+\varepsilon.
}
\]

This often corresponds more closely to numerical practice.

%%%%%%%%%%%%%%%%%%%%%%%%%%%%%%%%%%%%%%%%%%%%%%%%%%%%%%%%%%%%
\subsection{Error Bounds}
\label{subsec:constructive-error-bounds}
%%%%%%%%%%%%%%%%%%%%%%%%%%%%%%%%%%%%%%%%%%%%%%%%%%%%%%%%%%%%

An algorithm becomes mathematically useful when accompanied by a bound on
its error.

Constructive proofs often naturally produce statements such as

\[
\boxed{
|x_n-x|
\leq
E(n),
}
\]

where

\[
E(n)\to0.
\]

The function \(E\) gives explicit control over approximation quality.

%%%%%%%%%%%%%%%%%%%%%%%%%%%%%%%%%%%%%%%%%%%%%%%%%%%%%%%%%%%%
\subsection{Worked Example: Geometric Error Bound}
\label{subsec:worked-geometric-error-bound}
%%%%%%%%%%%%%%%%%%%%%%%%%%%%%%%%%%%%%%%%%%%%%%%%%%%%%%%%%%%%

\begin{workedexamplebox}{Choosing an Iteration Count}

Suppose an algorithm satisfies

\[
|x_n-x^*|
\leq
2^{-n}.
\]

To guarantee

\[
|x_n-x^*|
<
10^{-3},
\]

choose \(n\) such that

\[
2^{-n}
<
10^{-3}.
\]

Since

\[
2^{10}=1024,
\]

we have

\[
2^{-10}
<
10^{-3}.
\]

Thus

\begin{answerbox}
\[
\boxed{
n=10
}
\]

is sufficient.
\end{answerbox}

The convergence proof provides an explicit computational stopping rule.

\end{workedexamplebox}

%%%%%%%%%%%%%%%%%%%%%%%%%%%%%%%%%%%%%%%%%%%%%%%%%%%%%%%%%%%%
\subsection{Constructive Algebra}
\label{subsec:constructive-algebra}
%%%%%%%%%%%%%%%%%%%%%%%%%%%%%%%%%%%%%%%%%%%%%%%%%%%%%%%%%%%%

Many finite algebraic procedures are naturally constructive.

Examples include:

\[
\boxed{
\text{Euclidean algorithm},
}
\]

\[
\boxed{
\text{Gaussian elimination},
}
\]

and

\[
\boxed{
\text{finite polynomial computations}.
}
\]

The constructive content is especially clear when the algorithms terminate
with verifiable outputs.

%%%%%%%%%%%%%%%%%%%%%%%%%%%%%%%%%%%%%%%%%%%%%%%%%%%%%%%%%%%%
\subsection{Worked Example: Euclidean Algorithm}
\label{subsec:worked-euclidean-algorithm-constructive}
%%%%%%%%%%%%%%%%%%%%%%%%%%%%%%%%%%%%%%%%%%%%%%%%%%%%%%%%%%%%

\begin{workedexamplebox}{Constructing a Greatest Common Divisor}

Compute

\[
\gcd(252,105).
\]

Using the Euclidean algorithm,

\[
252
=
2(105)+42,
\]

\[
105
=
2(42)+21,
\]

\[
42
=
2(21)+0.
\]

Therefore,

\begin{answerbox}
\[
\boxed{
\gcd(252,105)=21.
}
\]
\end{answerbox}

The proof simultaneously gives an explicit algorithm for constructing the
greatest common divisor.

\end{workedexamplebox}

%%%%%%%%%%%%%%%%%%%%%%%%%%%%%%%%%%%%%%%%%%%%%%%%%%%%%%%%%%%%
\subsection{Constructive Linear Algebra}
\label{subsec:constructive-linear-algebra}
%%%%%%%%%%%%%%%%%%%%%%%%%%%%%%%%%%%%%%%%%%%%%%%%%%%%%%%%%%%%

Finite-dimensional linear algebra is often constructive when operations are
performed over fields with effective arithmetic.

Gaussian elimination can construct:

\[
\boxed{
\text{solutions},
}
\]

\[
\boxed{
\text{ranks},
}
\]

\[
\boxed{
\text{bases},
}
\]

and

\[
\boxed{
\text{determinants}.
}
\]

However, statements relying on arbitrary infinite-dimensional bases may
invoke stronger choice principles.

%%%%%%%%%%%%%%%%%%%%%%%%%%%%%%%%%%%%%%%%%%%%%%%%%%%%%%%%%%%%
\subsection{Choice and Constructive Mathematics}
\label{subsec:choice-constructive-mathematics}
%%%%%%%%%%%%%%%%%%%%%%%%%%%%%%%%%%%%%%%%%%%%%%%%%%%%%%%%%%%%

Choice principles behave differently in constructive foundations.

A proof of

\[
\forall x\in A
\exists y\in B
\,
R(x,y)
\]

may already contain a procedure selecting a witness \(y\) for each \(x\).

Thus some forms of choice become natural under propositions-as-types.

Other set-theoretic choice principles remain nonconstructive or require
additional axioms.

%%%%%%%%%%%%%%%%%%%%%%%%%%%%%%%%%%%%%%%%%%%%%%%%%%%%%%%%%%%%
\subsection{Countable Choice}
\label{subsec:countable-choice-constructive}
%%%%%%%%%%%%%%%%%%%%%%%%%%%%%%%%%%%%%%%%%%%%%%%%%%%%%%%%%%%%

Countable choice states, roughly, that from

\[
\forall n\in\mathbb{N}
\exists x\in A_n
\]

one may obtain a sequence

\[
\boxed{
(x_n)_{n\in\mathbb{N}}
}
\]

with

\[
x_n\in A_n.
\]

Its constructive status depends on the foundational system and the meaning
of the existential statements.

%%%%%%%%%%%%%%%%%%%%%%%%%%%%%%%%%%%%%%%%%%%%%%%%%%%%%%%%%%%%
\subsection{Dependent Choice}
\label{subsec:dependent-choice-constructive}
%%%%%%%%%%%%%%%%%%%%%%%%%%%%%%%%%%%%%%%%%%%%%%%%%%%%%%%%%%%%

Dependent choice permits recursively selecting a sequence when each next
choice depends on the previous one.

Given relation \(R\), one seeks

\[
x_0,x_1,x_2,\ldots
\]

such that

\[
\boxed{
R(x_n,x_{n+1})
}
\]

for every \(n\).

This principle is important in analysis and can be studied constructively
in various forms.

%%%%%%%%%%%%%%%%%%%%%%%%%%%%%%%%%%%%%%%%%%%%%%%%%%%%%%%%%%%%
\subsection{Constructive Set Theory}
\label{subsec:constructive-set-theory}
%%%%%%%%%%%%%%%%%%%%%%%%%%%%%%%%%%%%%%%%%%%%%%%%%%%%%%%%%%%%

Set theory can also be formulated constructively.

Such systems modify classical set theory by using intuitionistic logic and
carefully selecting axioms compatible with constructive interpretations.

The resulting theories preserve set-based foundations while avoiding
unrestricted classical reasoning.

%%%%%%%%%%%%%%%%%%%%%%%%%%%%%%%%%%%%%%%%%%%%%%%%%%%%%%%%%%%%
\subsection{Constructive Type Theory}
\label{subsec:constructive-type-theory}
%%%%%%%%%%%%%%%%%%%%%%%%%%%%%%%%%%%%%%%%%%%%%%%%%%%%%%%%%%%%

Dependent type theory provides another natural home for constructive
mathematics.

Under Curry--Howard:

\[
\boxed{
\exists x:A\,P(x)
}
\]

is represented by

\[
\boxed{
\sum_{x:A}P(x),
}
\]

and

\[
\boxed{
\forall x:A\,P(x)
}
\]

is represented by

\[
\boxed{
\prod_{x:A}P(x).
}
\]

Thus constructive evidence is built directly into the language.

%%%%%%%%%%%%%%%%%%%%%%%%%%%%%%%%%%%%%%%%%%%%%%%%%%%%%%%%%%%%
\subsection{Realizability}
\label{subsec:realizability-preview}
%%%%%%%%%%%%%%%%%%%%%%%%%%%%%%%%%%%%%%%%%%%%%%%%%%%%%%%%%%%%

Realizability interprets logical statements through computational objects
that realize them.

For example, a realizer of

\[
P\rightarrow Q
\]

acts as a computational procedure transforming realizers of \(P\) into
realizers of \(Q\).

This provides a formal bridge between logic and computation.

%%%%%%%%%%%%%%%%%%%%%%%%%%%%%%%%%%%%%%%%%%%%%%%%%%%%%%%%%%%%
\subsection{Realizers for Logical Connectives}
\label{subsec:realizers-logical-connectives}
%%%%%%%%%%%%%%%%%%%%%%%%%%%%%%%%%%%%%%%%%%%%%%%%%%%%%%%%%%%%

Conceptually:

a realizer of

\[
P\land Q
\]

contains realizers of both propositions;

a realizer of

\[
P\lor Q
\]

contains a tag and a realizer of the chosen side;

a realizer of

\[
\exists x\,P(x)
\]

contains a witness and a realizer of \(P(x)\).

This closely resembles Curry--Howard.

%%%%%%%%%%%%%%%%%%%%%%%%%%%%%%%%%%%%%%%%%%%%%%%%%%%%%%%%%%%%
\subsection{Proof Interpretations}
\label{subsec:constructive-proof-interpretations}
%%%%%%%%%%%%%%%%%%%%%%%%%%%%%%%%%%%%%%%%%%%%%%%%%%%%%%%%%%%%

Proof interpretations translate proofs into mathematical or computational
objects.

Their goals may include:

\[
\boxed{
\text{extracting algorithms},
}
\]

\[
\boxed{
\text{measuring proof strength},
}
\]

or

\[
\boxed{
\text{showing conservativity}.
}
\]

Constructive proof interpretations often reveal hidden computational content
inside classical arguments.

%%%%%%%%%%%%%%%%%%%%%%%%%%%%%%%%%%%%%%%%%%%%%%%%%%%%%%%%%%%%
\subsection{Double-Negation Translation}
\label{subsec:double-negation-translation}
%%%%%%%%%%%%%%%%%%%%%%%%%%%%%%%%%%%%%%%%%%%%%%%%%%%%%%%%%%%%

Classical reasoning can sometimes be embedded into intuitionistic logic by
systematically inserting double negations.

This shows that classical mathematics and constructive mathematics are not
completely disconnected.

Instead, classical proofs may sometimes be interpreted constructively in a
weaker, transformed form.

%%%%%%%%%%%%%%%%%%%%%%%%%%%%%%%%%%%%%%%%%%%%%%%%%%%%%%%%%%%%
\subsection{Constructive Content of Classical Proofs}
\label{subsec:constructive-content-classical-proofs}
%%%%%%%%%%%%%%%%%%%%%%%%%%%%%%%%%%%%%%%%%%%%%%%%%%%%%%%%%%%%

A classical theorem may contain useful computational information even if the
original proof is nonconstructive.

Proof transformations can sometimes extract:

\[
\boxed{
\text{bounds},
}
\]

\[
\boxed{
\text{approximation algorithms},
}
\]

or

\[
\boxed{
\text{finite witnesses}.
}
\]

This motivates proof mining and related areas.

%%%%%%%%%%%%%%%%%%%%%%%%%%%%%%%%%%%%%%%%%%%%%%%%%%%%%%%%%%%%
\subsection{Proof Mining Preview}
\label{subsec:proof-mining-preview}
%%%%%%%%%%%%%%%%%%%%%%%%%%%%%%%%%%%%%%%%%%%%%%%%%%%%%%%%%%%%

Proof mining analyzes ordinary mathematical proofs to extract explicit
quantitative information.

For example, a theorem asserting convergence may implicitly contain:

\[
\boxed{
\text{a rate of convergence},
}
\]

\[
\boxed{
\text{a metastability bound},
}
\]

or

\[
\boxed{
\text{an iteration count}.
}
\]

The goal is to uncover this hidden constructive content.

%%%%%%%%%%%%%%%%%%%%%%%%%%%%%%%%%%%%%%%%%%%%%%%%%%%%%%%%%%%%
\subsection{Constructive Numerical Mathematics}
\label{subsec:constructive-numerical-mathematics}
%%%%%%%%%%%%%%%%%%%%%%%%%%%%%%%%%%%%%%%%%%%%%%%%%%%%%%%%%%%%

Numerical mathematics is naturally connected with constructivism.

A numerical algorithm typically provides:

\[
\boxed{
\text{an explicit approximation},
}
\]

\[
\boxed{
\text{a stopping criterion},
}
\]

and

\[
\boxed{
\text{an error estimate}.
}
\]

These are exactly the kinds of data constructive mathematics values.

%%%%%%%%%%%%%%%%%%%%%%%%%%%%%%%%%%%%%%%%%%%%%%%%%%%%%%%%%%%%
\subsection{Constructive Root Finding}
\label{subsec:constructive-root-finding}
%%%%%%%%%%%%%%%%%%%%%%%%%%%%%%%%%%%%%%%%%%%%%%%%%%%%%%%%%%%%

Suppose an iterative procedure generates

\[
x_0,x_1,x_2,\ldots
\]

and one proves

\[
\boxed{
|x_n-r|
\leq
E(n),
}
\]

where

\[
E(n)\to0.
\]

Then for desired tolerance \(\varepsilon\), choose \(n\) such that

\[
E(n)<\varepsilon.
\]

The proof becomes an executable numerical guarantee.

%%%%%%%%%%%%%%%%%%%%%%%%%%%%%%%%%%%%%%%%%%%%%%%%%%%%%%%%%%%%
\subsection{Constructive Optimization}
\label{subsec:constructive-optimization}
%%%%%%%%%%%%%%%%%%%%%%%%%%%%%%%%%%%%%%%%%%%%%%%%%%%%%%%%%%%%

Optimization results can be stated constructively through algorithms and
error bounds.

Instead of merely asserting that a minimizer exists, one may construct a
sequence

\[
x_n
\]

satisfying

\[
\boxed{
f(x_n)-f^*
\leq
E(n),
}
\]

where

\[
E(n)\to0.
\]

The convergence estimate determines how many iterations are required to
reach a desired accuracy.

%%%%%%%%%%%%%%%%%%%%%%%%%%%%%%%%%%%%%%%%%%%%%%%%%%%%%%%%%%%%
\subsection{Constructive Fixed-Point Methods}
\label{subsec:constructive-fixed-point-methods}
%%%%%%%%%%%%%%%%%%%%%%%%%%%%%%%%%%%%%%%%%%%%%%%%%%%%%%%%%%%%

A contraction mapping provides a particularly constructive fixed-point
principle.

Suppose

\[
d(Tx,Ty)
\leq
q\,d(x,y),
\qquad
0\leq q<1.
\]

Starting with \(x_0\), define

\[
\boxed{
x_{n+1}=T(x_n).
}
\]

The iterates converge to the unique fixed point.

Moreover, explicit error estimates can be derived.

%%%%%%%%%%%%%%%%%%%%%%%%%%%%%%%%%%%%%%%%%%%%%%%%%%%%%%%%%%%%
\subsection{Worked Example: Contraction Error Bound}
\label{subsec:worked-contraction-error}
%%%%%%%%%%%%%%%%%%%%%%%%%%%%%%%%%%%%%%%%%%%%%%%%%%%%%%%%%%%%

\begin{workedexamplebox}{Constructive Fixed-Point Approximation}

Suppose

\[
d(Tx,Ty)
\leq
q\,d(x,y),
\qquad
0<q<1.
\]

Let

\[
x_{n+1}=T(x_n).
\]

A standard estimate gives

\[
d(x_n,x^*)
\leq
\frac{q^n}{1-q}
d(x_1,x_0).
\]

Thus to guarantee

\[
d(x_n,x^*)<\varepsilon,
\]

choose \(n\) satisfying

\[
\frac{q^n}{1-q}
d(x_1,x_0)
<
\varepsilon.
\]

\begin{answerbox}
\[
\boxed{
\text{The fixed-point theorem supplies both the solution and an
approximation algorithm.}
}
\]
\end{answerbox}

\end{workedexamplebox}

%%%%%%%%%%%%%%%%%%%%%%%%%%%%%%%%%%%%%%%%%%%%%%%%%%%%%%%%%%%%
\subsection{Constructive Probability}
\label{subsec:constructive-probability}
%%%%%%%%%%%%%%%%%%%%%%%%%%%%%%%%%%%%%%%%%%%%%%%%%%%%%%%%%%%%

Constructive probability requires careful representations of:

\[
\boxed{
\text{random variables},
}
\]

\[
\boxed{
\text{distributions},
}
\]

and

\[
\boxed{
\text{expectations}.
}
\]

Approximation methods such as finite distributions and computable sampling
procedures can provide constructive interpretations of probabilistic
objects.

%%%%%%%%%%%%%%%%%%%%%%%%%%%%%%%%%%%%%%%%%%%%%%%%%%%%%%%%%%%%
\subsection{Constructive Topology}
\label{subsec:constructive-topology}
%%%%%%%%%%%%%%%%%%%%%%%%%%%%%%%%%%%%%%%%%%%%%%%%%%%%%%%%%%%%

Constructive topology often emphasizes positive information.

Open sets are naturally suited to this viewpoint because membership can be
witnessed locally.

Classical principles involving arbitrary complements or arbitrary
subsets may require reformulation.

Locale and point-free approaches also interact with constructive
foundations.

%%%%%%%%%%%%%%%%%%%%%%%%%%%%%%%%%%%%%%%%%%%%%%%%%%%%%%%%%%%%
\subsection{Positive Information}
\label{subsec:positive-information-constructive}
%%%%%%%%%%%%%%%%%%%%%%%%%%%%%%%%%%%%%%%%%%%%%%%%%%%%%%%%%%%%

Constructive mathematics often distinguishes positive information from
negative information.

For example,

\[
\boxed{
x>0
}
\]

may contain a quantitative witness showing that \(x\) is bounded away from
zero.

By contrast,

\[
\boxed{
\neg(x\leq0)
}
\]

may not provide such a bound.

This distinction appears throughout constructive analysis.

%%%%%%%%%%%%%%%%%%%%%%%%%%%%%%%%%%%%%%%%%%%%%%%%%%%%%%%%%%%%
\subsection{Constructive Counterexamples}
\label{subsec:constructive-counterexamples}
%%%%%%%%%%%%%%%%%%%%%%%%%%%%%%%%%%%%%%%%%%%%%%%%%%%%%%%%%%%%

In constructive mathematics, failure to prove a classical principle does
not automatically mean its negation is provable.

A constructive counterexample may instead show that accepting a principle
would imply another nonconstructive principle.

Thus independence and implication relationships are important.

%%%%%%%%%%%%%%%%%%%%%%%%%%%%%%%%%%%%%%%%%%%%%%%%%%%%%%%%%%%%
\subsection{Classical Theorems With Constructive Versions}
\label{subsec:classical-theorems-constructive-versions}
%%%%%%%%%%%%%%%%%%%%%%%%%%%%%%%%%%%%%%%%%%%%%%%%%%%%%%%%%%%%

Many classical theorems have constructive counterparts after modifying the
statement.

Common modifications include replacing:

\[
\boxed{
\text{exact existence}
}
\]

with

\[
\boxed{
\text{arbitrarily accurate approximation},
}
\]

or adding:

\[
\boxed{
\text{locatedness},
}
\]

\[
\boxed{
\text{uniform continuity},
}
\]

or

\[
\boxed{
\text{effective compactness}.
}
\]

%%%%%%%%%%%%%%%%%%%%%%%%%%%%%%%%%%%%%%%%%%%%%%%%%%%%%%%%%%%%
\subsection{Constructive Strengthening}
\label{subsec:constructive-strengthening}
%%%%%%%%%%%%%%%%%%%%%%%%%%%%%%%%%%%%%%%%%%%%%%%%%%%%%%%%%%%%

A constructive theorem can sometimes be more informative than its classical
counterpart.

For example, instead of merely proving

\[
\exists x\,P(x),
\]

one may construct an algorithm

\[
\boxed{
x=F(\text{input data})
}
\]

and derive an explicit complexity or error bound.

Thus constructive restrictions can lead to stronger quantitative
conclusions.

%%%%%%%%%%%%%%%%%%%%%%%%%%%%%%%%%%%%%%%%%%%%%%%%%%%%%%%%%%%%
\subsection{Constructivism and Formal Verification}
\label{subsec:constructivism-formal-verification}
%%%%%%%%%%%%%%%%%%%%%%%%%%%%%%%%%%%%%%%%%%%%%%%%%%%%%%%%%%%%

Formal verification benefits from constructive proof because proof terms can
carry executable evidence.

A machine-checkable theorem of the form

\[
\forall x
\exists y\,R(x,y)
\]

may correspond directly to verified code computing \(y\) from \(x\).

This makes constructive foundations especially attractive for proof
assistants.

%%%%%%%%%%%%%%%%%%%%%%%%%%%%%%%%%%%%%%%%%%%%%%%%%%%%%%%%%%%%
\subsection{Constructive Mathematics and Type Theory}
\label{subsec:constructive-mathematics-type-theory}
%%%%%%%%%%%%%%%%%%%%%%%%%%%%%%%%%%%%%%%%%%%%%%%%%%%%%%%%%%%%

Type theory internalizes the constructive meaning of propositions.

Under propositions-as-types:

\[
\boxed{
P\lor Q
}
\]

becomes a sum type,

\[
\boxed{
P\land Q
}
\]

becomes a product type,

and

\[
\boxed{
\exists x:A\,P(x)
}
\]

becomes a dependent pair.

Thus constructive logic and type theory fit together naturally.

%%%%%%%%%%%%%%%%%%%%%%%%%%%%%%%%%%%%%%%%%%%%%%%%%%%%%%%%%%%%
\subsection{Constructive Mathematics and Computability}
\label{subsec:constructive-mathematics-computability}
%%%%%%%%%%%%%%%%%%%%%%%%%%%%%%%%%%%%%%%%%%%%%%%%%%%%%%%%%%%%

Computability theory supplies formal notions of algorithm and effective
procedure.

Constructive mathematics uses related ideas inside mathematical proofs.

However, not every constructive system identifies ``constructive'' with
``Turing computable'' in exactly the same way.

The foundational framework determines the precise interpretation.

%%%%%%%%%%%%%%%%%%%%%%%%%%%%%%%%%%%%%%%%%%%%%%%%%%%%%%%%%%%%
\subsection{Constructive Mathematics and Proof Theory}
\label{subsec:constructive-mathematics-proof-theory}
%%%%%%%%%%%%%%%%%%%%%%%%%%%%%%%%%%%%%%%%%%%%%%%%%%%%%%%%%%%%

Proof theory studies transformations of constructive proofs.

Normalization can expose computational content.

Translations can compare classical and intuitionistic systems.

Ordinal and proof-theoretic methods can measure the strength of constructive
formal theories.

%%%%%%%%%%%%%%%%%%%%%%%%%%%%%%%%%%%%%%%%%%%%%%%%%%%%%%%%%%%%
\subsection{Constructive Mathematics Workflow}
\label{subsec:constructive-mathematics-workflow}
%%%%%%%%%%%%%%%%%%%%%%%%%%%%%%%%%%%%%%%%%%%%%%%%%%%%%%%%%%%%

When seeking a constructive proof:

\begin{enumerate}
    \item identify the logical form of the target statement;
    \item if the goal is existential, determine what witness must be built;
    \item if the goal is a disjunction, determine which side can be proved;
    \item if the goal is an implication, construct a transformation from
          hypotheses to conclusion;
    \item if the goal is universal, construct a uniform method for arbitrary
          inputs;
    \item distinguish positive evidence from mere failure of contradiction;
    \item avoid unrestricted use of excluded middle unless the proposition
          is decidable or the chosen system allows it;
    \item avoid unrestricted double-negation elimination;
    \item use contradiction freely when proving a negation;
    \item identify finite searches and perform them explicitly;
    \item supply convergence rates when using sequences;
    \item supply moduli when using continuity or Cauchy convergence;
    \item use apartness when positive inequality information is needed;
    \item formulate approximate versions when exact existence is too strong;
    \item include quantitative error bounds;
    \item identify whether choice principles are being used;
    \item convert proofs into algorithms when possible;
    \item verify termination of recursive constructions;
    \item distinguish constructivity from computational efficiency;
    \item state clearly which constructive foundation is being assumed.
\end{enumerate}

%%%%%%%%%%%%%%%%%%%%%%%%%%%%%%%%%%%%%%%%%%%%%%%%%%%%%%%%%%%%
\subsection{Common Errors}
\label{subsec:constructive-mathematics-common-errors}
%%%%%%%%%%%%%%%%%%%%%%%%%%%%%%%%%%%%%%%%%%%%%%%%%%%%%%%%%%%%

\subsubsection{Assuming Constructive Mathematics Rejects Logic}

Constructive mathematics uses rigorous logic.

It generally uses intuitionistic rather than unrestricted classical logic.

%%%%%%%%%%%%%%%%%%%%%%%%%%%%%%%%%%%%%%%%%%%%%%%%%%%%%%%%%%%%
\subsubsection{Assuming Constructive Mathematics Rejects All Proofs by Contradiction}

Proof by contradiction is constructive when proving a negation.

The difficulty arises when one proves only

\[
\neg\neg P
\]

and then concludes arbitrary positive \(P\).

%%%%%%%%%%%%%%%%%%%%%%%%%%%%%%%%%%%%%%%%%%%%%%%%%%%%%%%%%%%%
\subsubsection{Assuming Excluded Middle Is Always False Constructively}

Constructive logic does not prove

\[
P\lor\neg P
\]

for every arbitrary proposition.

For decidable propositions, excluded middle can be proved.

%%%%%%%%%%%%%%%%%%%%%%%%%%%%%%%%%%%%%%%%%%%%%%%%%%%%%%%%%%%%
\subsubsection{Confusing \(\neg\neg P\) With \(P\)}

Classically these are equivalent.

Constructively,

\[
\boxed{
P\rightarrow\neg\neg P
}
\]

holds, but

\[
\boxed{
\neg\neg P\rightarrow P
}
\]

does not hold universally.

%%%%%%%%%%%%%%%%%%%%%%%%%%%%%%%%%%%%%%%%%%%%%%%%%%%%%%%%%%%%
\subsubsection{Treating an Existential Contradiction Proof as an Explicit Witness}

A classical proof that

\[
\neg
\forall x\,\neg P(x)
\]

does not automatically produce an \(x\) satisfying \(P(x)\).

Constructive existence requires witness information.

%%%%%%%%%%%%%%%%%%%%%%%%%%%%%%%%%%%%%%%%%%%%%%%%%%%%%%%%%%%%
\subsubsection{Assuming Constructive Means Numerically Approximate}

Constructive mathematics includes exact finite and symbolic constructions.

Approximation becomes important especially for infinite or continuous
objects.

%%%%%%%%%%%%%%%%%%%%%%%%%%%%%%%%%%%%%%%%%%%%%%%%%%%%%%%%%%%%
\subsubsection{Assuming Every Algorithm Is Efficient}

Constructivity concerns obtaining an effective procedure.

The procedure may still be computationally expensive.

Complexity analysis is a separate question.

%%%%%%%%%%%%%%%%%%%%%%%%%%%%%%%%%%%%%%%%%%%%%%%%%%%%%%%%%%%%
\subsubsection{Ignoring Error Bounds}

Providing numerical approximations without knowing their accuracy may not
supply enough constructive information.

A constructive algorithm should ideally include a way to control error.

%%%%%%%%%%%%%%%%%%%%%%%%%%%%%%%%%%%%%%%%%%%%%%%%%%%%%%%%%%%%
\subsubsection{Ignoring Convergence Rates}

A sequence may converge classically without an available effective rate of
convergence.

Constructive applications often require quantitative convergence data.

%%%%%%%%%%%%%%%%%%%%%%%%%%%%%%%%%%%%%%%%%%%%%%%%%%%%%%%%%%%%
\subsubsection{Confusing \(x\neq y\) With Apartness}

The statement

\[
x\neq y
\]

means equality gives contradiction.

Apartness

\[
x\# y
\]

contains positive evidence separating the two objects.

%%%%%%%%%%%%%%%%%%%%%%%%%%%%%%%%%%%%%%%%%%%%%%%%%%%%%%%%%%%%
\subsubsection{Assuming Equality of Arbitrary Reals Is Decidable}

Constructive real equality is generally not decidable.

Finite approximations may never determine exact equality.

%%%%%%%%%%%%%%%%%%%%%%%%%%%%%%%%%%%%%%%%%%%%%%%%%%%%%%%%%%%%
\subsubsection{Assuming Every Classical Compactness Equivalence Remains Constructively Valid}

Classical topology contains many equivalent compactness properties.

Constructively, these may separate and require distinct hypotheses.

%%%%%%%%%%%%%%%%%%%%%%%%%%%%%%%%%%%%%%%%%%%%%%%%%%%%%%%%%%%%
\subsubsection{Assuming Every Classical Existence Theorem Has the Same Constructive Statement}

Some classical theorems require reformulation through approximation,
locatedness, explicit continuity, or additional data.

%%%%%%%%%%%%%%%%%%%%%%%%%%%%%%%%%%%%%%%%%%%%%%%%%%%%%%%%%%%%
\subsubsection{Confusing Constructive Choice With Classical Choice}

Some forms of choice follow naturally when existential proofs contain
witnesses.

Other forms remain independent or nonconstructive depending on the
foundation.

%%%%%%%%%%%%%%%%%%%%%%%%%%%%%%%%%%%%%%%%%%%%%%%%%%%%%%%%%%%%
\subsubsection{Assuming Constructive Mathematics Is One Formal Theory}

There are several constructive frameworks.

Bishop-style mathematics, intuitionistic set theory, constructive set
theory, and dependent type theory may differ in accepted principles.

%%%%%%%%%%%%%%%%%%%%%%%%%%%%%%%%%%%%%%%%%%%%%%%%%%%%%%%%%%%%
\subsubsection{Confusing Constructivism With Finitism}

Constructive mathematics can reason about infinite objects.

It requires appropriate constructive representations rather than restricting
all mathematics to finite objects.

%%%%%%%%%%%%%%%%%%%%%%%%%%%%%%%%%%%%%%%%%%%%%%%%%%%%%%%%%%%%
\subsubsection{Confusing Constructivism With Predicativism}

Constructivism concerns proof and existence principles.

Predicativism concerns restrictions on definitions involving totalities that
include the object being defined.

The ideas overlap historically but are not identical.

%%%%%%%%%%%%%%%%%%%%%%%%%%%%%%%%%%%%%%%%%%%%%%%%%%%%%%%%%%%%
\subsubsection{Assuming Nonconstructive Proofs Are Invalid Classically}

Classical mathematics accepts many nonconstructive arguments.

Constructive mathematics uses a stronger information standard, not a claim
that classical deductions are syntactically invalid in classical logic.

%%%%%%%%%%%%%%%%%%%%%%%%%%%%%%%%%%%%%%%%%%%%%%%%%%%%%%%%%%%%
\subsection{Applications}
\label{subsec:constructive-mathematics-applications}
%%%%%%%%%%%%%%%%%%%%%%%%%%%%%%%%%%%%%%%%%%%%%%%%%%%%%%%%%%%%

\begin{applicationbox}

\textbf{Numerical Analysis.}
Constructive proofs naturally produce algorithms, convergence rates,
stopping criteria, and error bounds.

\medskip

\textbf{Optimization.}
Approximate minimizers and explicit convergence guarantees provide
constructive versions of existence results.

\medskip

\textbf{Computer-Assisted Mathematics.}
Explicit witnesses and finite verification procedures fit naturally with
machine-checked mathematics.

\medskip

\textbf{Formal Verification.}
Constructive existence proofs can contain executable programs satisfying
formal specifications.

\medskip

\textbf{Type Theory.}
The propositions-as-types interpretation gives a direct formal realization
of constructive logic.

\medskip

\textbf{Proof Theory.}
Normalization and proof interpretation can extract computational content
from derivations.

\medskip

\textbf{Computability Theory.}
Effective procedures provide formal models for many constructive
constructions.

\medskip

\textbf{Analysis.}
Constructive real numbers, moduli of continuity, effective compactness, and
quantitative approximation create computational versions of analytic
theorems.

\medskip

\textbf{Foundations of Mathematics.}
Constructive mathematics provides an alternative understanding of
existence, truth, and mathematical proof.

\end{applicationbox}

%%%%%%%%%%%%%%%%%%%%%%%%%%%%%%%%%%%%%%%%%%%%%%%%%%%%%%%%%%%%
\subsection{Exercises}
\label{subsec:constructive-mathematics-exercises}
%%%%%%%%%%%%%%%%%%%%%%%%%%%%%%%%%%%%%%%%%%%%%%%%%%%%%%%%%%%%

\begin{exercise}[1]
Define constructive existence.
\end{exercise}

\begin{exercise}[1]
State the law of excluded middle.
\end{exercise}

\begin{exercise}[1]
State double-negation elimination.
\end{exercise}

\begin{exercise}[1]
Define a decidable proposition.
\end{exercise}

\begin{exercise}[1]
Define apartness conceptually.
\end{exercise}

\begin{exercise}[1]
Define a modulus of convergence.
\end{exercise}

\begin{exercise}[1]
Define a modulus of continuity.
\end{exercise}

\begin{exercise}[1]
Define total boundedness.
\end{exercise}

\begin{exercise}[1]
Explain realizability conceptually.
\end{exercise}

\begin{exercise}[1]
Explain algorithm extraction conceptually.
\end{exercise}

\begin{exercise}[2]
Give a constructive proof that there exists an even natural number greater
than \(10\).
\end{exercise}

\begin{exercise}[2]
Give an explicit witness for

\[
\exists n\in\mathbb{N}
\,
(n^2=49).
\]
\end{exercise}

\begin{exercise}[2]
Explain why a constructive proof of

\[
P\lor Q
\]

must identify a side.
\end{exercise}

\begin{exercise}[2]
Explain why

\[
P\rightarrow\neg\neg P
\]

is constructively valid.
\end{exercise}

\begin{exercise}[3]
Explain why

\[
\neg\neg P\rightarrow P
\]

is not intuitionistically valid in general.
\end{exercise}

\begin{exercise}[2]
Explain why proof by contradiction is constructive when the goal is
\(\neg P\).
\end{exercise}

\begin{exercise}[3]
Rewrite a classical contradiction proof so that you identify exactly where
double-negation elimination is used.
\end{exercise}

\begin{exercise}[2]
Show that equality of two given natural numbers is decidable.
\end{exercise}

\begin{exercise}[2]
Show that divisibility of two given integers is decidable.
\end{exercise}

\begin{exercise}[2]
Explain why finite search is constructive.
\end{exercise}

\begin{exercise}[3]
Given a finite list and a decidable property, write a constructive procedure
for finding a witness or proving no witness exists.
\end{exercise}

\begin{exercise}[3]
Explain the BHK interpretation of conjunction.
\end{exercise}

\begin{exercise}[3]
Explain the BHK interpretation of disjunction.
\end{exercise}

\begin{exercise}[3]
Explain the BHK interpretation of implication.
\end{exercise}

\begin{exercise}[3]
Explain the BHK interpretation of existential quantification.
\end{exercise}

\begin{exercise}[3]
Explain the BHK interpretation of universal quantification.
\end{exercise}

\begin{exercise}[3]
Define stability of a proposition.
\end{exercise}

\begin{exercise}[3]
Show that propositions of the form \(\neg P\) are stable intuitionistically.
\end{exercise}

\begin{exercise}[3]
Explain why Markov-style principles go beyond pure intuitionistic logic.
\end{exercise}

\begin{exercise}[2]
Explain why natural-number induction is constructive.
\end{exercise}

\begin{exercise}[2]
Construct a witness function for

\[
\forall n\in\mathbb{N}
\exists m\in\mathbb{N}
\,
(m=n+2).
\]
\end{exercise}

\begin{exercise}[3]
Show how a proof of

\[
\forall x\exists y\,R(x,y)
\]

can encode a function \(f(x)=y\) in a propositions-as-types setting.
\end{exercise}

\begin{exercise}[2]
Describe a constructive representation of a real number using rational
approximations.
\end{exercise}

\begin{exercise}[3]
Explain why convergence information must accompany a constructive Cauchy
sequence.
\end{exercise}

\begin{exercise}[3]
Find an explicit \(N(\varepsilon)\) such that

\[
\left|
\frac{1}{n+1}
\right|
<
\varepsilon
\]

for all \(n\geq N(\varepsilon)\).
\end{exercise}

\begin{exercise}[3]
Find a modulus of convergence for

\[
q_n=2^{-n}.
\]
\end{exercise}

\begin{exercise}[3]
Explain why equality of arbitrary real numbers is not generally decidable
constructively.
\end{exercise}

\begin{exercise}[3]
Explain why apartness gives stronger information than inequality.
\end{exercise}

\begin{exercise}[3]
Give a positive witness that

\[
2\#3
\]

under the usual real-number apartness interpretation.
\end{exercise}

\begin{exercise}[3]
Explain locatedness conceptually.
\end{exercise}

\begin{exercise}[2]
State uniform continuity.
\end{exercise}

\begin{exercise}[3]
Find a modulus of continuity for

\[
f(x)=2x
\]

on \(\mathbb{R}\).
\end{exercise}

\begin{exercise}[3]
Find a modulus of uniform continuity for

\[
f(x)=x^2
\]

on \([0,1]\).
\end{exercise}

\begin{exercise}[3]
Explain why constructive root finding may emphasize approximate roots.
\end{exercise}

\begin{exercise}[3]
For bisection on an interval of initial width \(1\), determine how many
steps guarantee interval width below \(10^{-6}\).
\end{exercise}

\begin{exercise}[2]
Define an \(\varepsilon\)-net.
\end{exercise}

\begin{exercise}[3]
Construct an \(\varepsilon\)-net for \([0,1]\).
\end{exercise}

\begin{exercise}[3]
Explain why total boundedness contains computational information.
\end{exercise}

\begin{exercise}[3]
Explain the constructive relationship between completeness and Cauchy data.
\end{exercise}

\begin{exercise}[3]
Explain why exact maximum attainment may be more difficult constructively
than approximate maximization.
\end{exercise}

\begin{exercise}[3]
Formulate an \(\varepsilon\)-optimal version of a minimization problem.
\end{exercise}

\begin{exercise}[3]
Given

\[
f(x_n)-f^*
\leq
\frac1n,
\]

find a sufficient \(n\) to guarantee error below \(10^{-4}\).
\end{exercise}

\begin{exercise}[3]
Given

\[
|x_n-x^*|
\leq
3^{-n},
\]

find a sufficient \(n\) to guarantee error below \(10^{-5}\).
\end{exercise}

\begin{exercise}[2]
Explain why the Euclidean algorithm is constructive.
\end{exercise}

\begin{exercise}[2]
Use the Euclidean algorithm to construct

\[
\gcd(391,299).
\]
\end{exercise}

\begin{exercise}[3]
Explain why Gaussian elimination is naturally constructive in finite
dimensions.
\end{exercise}

\begin{exercise}[3]
Explain why arbitrary bases of infinite-dimensional vector spaces may
involve Choice.
\end{exercise}

\begin{exercise}[3]
Explain countable choice conceptually.
\end{exercise}

\begin{exercise}[3]
Explain dependent choice conceptually.
\end{exercise}

\begin{exercise}[3]
Explain why choice principles behave differently in type theory and set
theory.
\end{exercise}

\begin{exercise}[3]
Describe constructive set theory conceptually.
\end{exercise}

\begin{exercise}[3]
Explain how dependent type theory supports constructive mathematics.
\end{exercise}

\begin{exercise}[3]
Explain realizability for

\[
P\rightarrow Q.
\]
\end{exercise}

\begin{exercise}[3]
Explain realizability for

\[
P\lor Q.
\]
\end{exercise}

\begin{exercise}[3]
Explain realizability for

\[
\exists x\,P(x).
\]
\end{exercise}

\begin{exercise}[3]
Explain a double-negation translation conceptually.
\end{exercise}

\begin{exercise}[3]
Explain what proof mining attempts to extract.
\end{exercise}

\begin{exercise}[3]
Give an example of quantitative information that may be hidden in a
qualitative convergence proof.
\end{exercise}

\begin{exercise}[3]
Explain why numerical stopping criteria are constructive data.
\end{exercise}

\begin{exercise}[3]
For a contraction constant \(q=\frac12\), derive an iteration count sufficient
to make

\[
\frac{q^n}{1-q}d(x_1,x_0)
<
\varepsilon.
\]
\end{exercise}

\begin{exercise}[3]
Explain why a contraction mapping theorem has strong constructive content.
\end{exercise}

\begin{exercise}[3]
Distinguish constructive mathematics from computability theory.
\end{exercise}

\begin{exercise}[3]
Distinguish constructivism from finitism.
\end{exercise}

\begin{exercise}[3]
Distinguish constructivism from predicativism.
\end{exercise}

\begin{exercise}[3]
Explain why a classical theorem may still have computational content.
\end{exercise}

\begin{exercise}[3]
Explain how proof assistants benefit from constructive proof terms.
\end{exercise}

\begin{exercise}[4]
Study intuitionistic propositional logic formally.
\end{exercise}

\begin{exercise}[4]
Study intuitionistic first-order logic.
\end{exercise}

\begin{exercise}[4]
Study Kripke semantics for intuitionistic logic.
\end{exercise}

\begin{exercise}[4]
Study Heyting algebras.
\end{exercise}

\begin{exercise}[4]
Study the double-negation translation of classical logic.
\end{exercise}

\begin{exercise}[4]
Study realizability interpretations.
\end{exercise}

\begin{exercise}[4]
Study recursive realizability.
\end{exercise}

\begin{exercise}[4]
Study Bishop-style constructive analysis.
\end{exercise}

\begin{exercise}[4]
Study constructive real-number representations.
\end{exercise}

\begin{exercise}[4]
Study constructive compactness in metric spaces.
\end{exercise}

\begin{exercise}[4]
Study constructive versions of the Intermediate Value Theorem.
\end{exercise}

\begin{exercise}[4]
Study constructive versions of fixed-point theorems.
\end{exercise}

\begin{exercise}[4]
Study constructive measure and integration.
\end{exercise}

\begin{exercise}[4]
Study constructive functional analysis.
\end{exercise}

\begin{exercise}[4]
Study constructive set theories.
\end{exercise}

\begin{exercise}[4]
Study the constructive status of countable choice.
\end{exercise}

\begin{exercise}[4]
Study dependent choice constructively.
\end{exercise}

\begin{exercise}[4]
Study program extraction from formal proofs.
\end{exercise}

\begin{exercise}[4]
Study proof mining and quantitative interpretations.
\end{exercise}

\begin{exercise}[4]
Compare Bishop-style mathematics with dependent type theory.
\end{exercise}

%%%%%%%%%%%%%%%%%%%%%%%%%%%%%%%%%%%%%%%%%%%%%%%%%%%%%%%%%%%%
\subsection{Conceptual Summary}
\label{subsec:constructive-mathematics-summary}
%%%%%%%%%%%%%%%%%%%%%%%%%%%%%%%%%%%%%%%%%%%%%%%%%%%%%%%%%%%%

Constructive mathematics interprets proof as explicit mathematical evidence.

A proof of

\[
\boxed{
\exists x\,P(x)
}
\]

should provide a witness.

A proof of

\[
\boxed{
P\lor Q
}
\]

should identify a valid side.

A proof of

\[
\boxed{
P\rightarrow Q
}
\]

should transform evidence for \(P\) into evidence for \(Q\).

Constructive logic therefore does not accept unrestricted use of

\[
\boxed{
P\lor\neg P
}
\]

or

\[
\boxed{
\neg\neg P\rightarrow P.
}
\]

Natural numbers, finite search, recursion, and induction are highly
constructive.

For real and infinite objects, constructive mathematics emphasizes:

\[
\boxed{
\text{explicit approximations},
}
\]

\[
\boxed{
\text{rates of convergence},
}
\]

\[
\boxed{
\text{moduli of continuity},
}
\]

\[
\boxed{
\text{apartness},
}
\]

and

\[
\boxed{
\text{quantitative error bounds}.
}
\]

Constructive mathematics therefore connects mathematical existence with
computation.

Through type theory and Curry--Howard,

\[
\boxed{
\text{proofs}
\leftrightarrow
\text{programs},
}
\]

while numerical and computational mathematics provide practical examples
of constructive procedures.

The constructive viewpoint does not merely ask whether a theorem is true.

It asks:

\[
\boxed{
\text{What explicit mathematical information does the proof provide?}
}
\]

%%%%%%%%%%%%%%%%%%%%%%%%%%%%%%%%%%%%%%%%%%%%%%%%%%%%%%%%%%%%
\subsection{Section Connections}
\label{subsec:constructive-mathematics-connections}
%%%%%%%%%%%%%%%%%%%%%%%%%%%%%%%%%%%%%%%%%%%%%%%%%%%%%%%%%%%%

\chapterconnection{
Constructive Mathematics develops the computational interpretation of proof
that appeared throughout Proof Theory, Computability Theory, Recursion
Theory, and especially Type Theory. Whereas classical mathematics often
accepts indirect existence arguments, constructive mathematics emphasizes
explicit witnesses, algorithms, approximation procedures, moduli, and error
bounds. These ideas connect foundational logic directly with analysis,
numerical mathematics, optimization, and formal verification. The final
section of Chapter~12, Philosophy and Foundations of Mathematics, compares
the broader foundational viewpoints underlying classical set theory,
formalism, intuitionism, constructivism, logicism, structuralism, and other
approaches to the nature of mathematical objects and mathematical truth.
}

%%%%%%%%%%%%%%%%%%%%%%%%%%%%%%%%%%%%%%%%%%%%%%%%%%%%%%%%%%%%
% SECTION SOURCE TRACKING
%%%%%%%%%%%%%%%%%%%%%%%%%%%%%%%%%%%%%%%%%%%%%%%%%%%%%%%%%%%%

%------------------------------------------------------------
% 12.8 Constructive Mathematics
%------------------------------------------------------------
%
% Source basis:
%   - Standard intuitionistic and constructive mathematics.
%   - Standard Bishop-style constructive analysis.
%   - Standard propositions-as-types and realizability viewpoints.
%
% Used for:
%   - constructive existence;
%   - explicit witnesses;
%   - classical versus constructive proof;
%   - intuitionistic logic;
%   - excluded middle;
%   - decidable propositions;
%   - finite search;
%   - double negation;
%   - constructive proof by contradiction;
%   - proof of negation;
%   - constructive disjunction;
%   - constructive implication;
%   - universal and existential quantification;
%   - BHK interpretation;
%   - constructive truth;
%   - stability;
%   - negative formulas;
%   - Markov-type principles preview;
%   - varieties of constructivism;
%   - Bishop-style mathematics;
%   - constructive versus computable mathematics;
%   - constructive natural numbers;
%   - constructive induction and recursion;
%   - algorithm extraction;
%   - constructive real numbers;
%   - Cauchy representations;
%   - convergence moduli;
%   - effective approximation;
%   - real equality;
%   - apartness;
%   - positive order;
%   - locatedness;
%   - constructive suprema;
%   - constructive continuity;
%   - uniform continuity;
%   - continuity moduli;
%   - intermediate-value issues;
%   - approximate roots;
%   - constructive compactness;
%   - total boundedness;
%   - completeness;
%   - epsilon-nets;
%   - approximate optimization;
%   - exact versus approximate solutions;
%   - error bounds;
%   - constructive algebra;
%   - constructive linear algebra;
%   - choice principles;
%   - countable choice;
%   - dependent choice;
%   - constructive set theory;
%   - constructive type theory;
%   - realizability;
%   - proof interpretations;
%   - double-negation translations;
%   - proof mining;
%   - constructive numerical mathematics;
%   - fixed-point approximation;
%   - constructive probability;
%   - constructive topology;
%   - positive information;
%   - constructive counterexamples;
%   - formal verification;
%   - applications and exercises.
%
% Notes:
%   - Philosophy and Foundations of Mathematics is developed next in
%     Section 12.9.
%   - Proof Theory appears in Section 12.4.
%   - Computability Theory appears in Section 12.5.
%   - Recursion Theory appears in Section 12.6.
%   - Type Theory appears in Section 12.7.
%
%%%%%%%%%%%%%%%%%%%%%%%%%%%%%%%%%%%%%%%%%%%%%%%%%%%%%%%%%%%%